\documentclass[11pt]{article}
\usepackage[margin=1in]{geometry}
\usepackage{amsmath,amssymb,amsthm}
\usepackage[T1]{fontenc}
\usepackage{lmodern}
\usepackage[colorlinks=true,linkcolor=blue,urlcolor=blue,citecolor=blue]{hyperref}
\newtheorem{theorem}{Theorem}
\newtheorem{lemma}{Lemma}
\newcommand{\T}{\mathbb T}
\newcommand{\R}{\mathbb R}
\newcommand{\Z}{\mathbb Z}
\newcommand{\dd}{\mathrm d}
\title{A positive-volume ugly band in a smooth Hamiltonian system}
\author{Alec Kriebel\\\small Independent researcher\\
\small\href{https://orcid.org/0009-0001-9320-500X}{ORCID: 0009-0001-9320-500X}}
\date{October 12, 2026\quad(version 1.0.0)}
\begin{document}
\maketitle
\begin{abstract}
We construct a compact connected symplectic four-manifold and a smooth,
real-valued autonomous Hamiltonian with an invariant open band of normalized
Liouville volume $1/4$. Every Lyapunov exponent at every point of the band is
zero, but no invariant measurable subset of positive volume supports an
almost-periodic Koopman representation. Each energy component in the band
is itself ergodic and non-almost-periodic. This gives an affirmative answer to
the existential question in item 10 of Knill's 2000 Hamiltonian problem list,
in the unrestricted smooth symplectic category and with almost periodicity
defined by precompact $L^2$ orbits. The proof combines the classical irrational
skew shift with a symplectic mapping torus and an explicit smooth cutoff.
Historical first priority is not asserted. AI tools were used extensively;
this research note is unrefereed and no human peer review is claimed.
\end{abstract}

\section{Question, convention, and result}

Item 10 of Knill's problem list \cite{knill2000} asks whether a Hamiltonian
system can have positive Liouville volume outside both its maximal
almost-periodic invariant region and its positive-Lyapunov-exponent region.
We call this complement the \emph{ugly set}. The printed existential question
does not specify a Euclidean phase space or an exact symplectic form.
We give an example on a compact smooth symplectic manifold. The construction
does not address additional requirements of analyticity, natural mechanics,
contact-type energies, genericity, or a prescribed cotangent bundle.

For a finite measure-preserving flow $\phi_s$ on $(X,\mu)$, we say that
the system is \emph{almost periodic} if every orbit
$\{f\circ\phi_s:s\in\R\}$ is relatively compact in $L^2(X,\mu)$.
An almost-periodic \emph{region} is an invariant measurable subset equipped
with its restricted measure; an almost-periodic factor or individual
eigenfunction is insufficient. This spectral convention is consistent with
Knill's later explicit terminology \cite[Section 4]{knill2016}; it is a stated
interpretation, since the 2000 item does not give a formal definition.
The proof below establishes noncompactness directly.

\begin{theorem}\label{thm:main}
There are a compact connected smooth symplectic four-manifold $(M,\omega)$,
a globally real-valued $H\in C^\infty(M)$, and an invariant open set $B$
with the following properties. The Liouville measure
$\mu=\omega^2/2$ has $\mu(M)=1$ and $\mu(B)=1/4$.
Every Lyapunov exponent of the Hamiltonian flow at every point of $B$ is zero.
No invariant measurable $C\subset B$ with $\mu(C)>0$ is almost periodic.
Each connected energy component in $B$, with its normalized conditional
Liouville measure, is ergodic and non-almost-periodic.
Consequently the ugly set has measure at least $1/4$.
\end{theorem}

We use $\iota_{X_H}\omega=\dd H$. Reversing the convention reverses
time and changes none of these assertions. The measure of a region always
means full four-dimensional Liouville measure unless explicitly described
as a conditional energy measure.

\section{The skew shift and its suspension}

Put $a=\sqrt2$ and let $m_2$ be normalized Haar measure on
$\T^2=\R^2/\Z^2$. The classical skew shift is
\[
 F(x,y)=(x+a,y+x)\pmod{\Z^2}.
\]
It is a smooth area-preserving diffeomorphism: integer changes of lifts
give integer changes of images, its inverse is
$(x-a,y-x+a)$, and its derivative is
$A=\left(\begin{smallmatrix}1&0\\1&1\end{smallmatrix}\right)$.
For every $n\in\Z$,
\begin{equation}\label{eq:iterate}
 F^n(x,y)=\left(x+na,y+nx+\frac{a n(n-1)}2\right),\qquad
 D F^n=\begin{pmatrix}1&0\\n&1\end{pmatrix}.
\end{equation}
These formulas follow by forward and backward induction.

The Fourier characters $e_{k,l}(x,y)=\exp(2\pi i(kx+ly))$
form an orthonormal basis, and
\[
 e_{k,l}\circ F=\exp(2\pi i k a)\,e_{k+l,l}.
\]
For an $F$-invariant $L^2$ function, Fourier coefficients have equal
absolute values along each chain $(k+nl,l)$, $n\in\Z$.
If $l\ne0$, square summability forces these coefficients to vanish.
If $l=0$, irrationality of $a$ eliminates every coefficient with $k\ne0$.
Thus $F$ is ergodic. Independently, $g(x,y)=\exp(2\pi i y)$ has
orthonormal integer translates under $F$, since for $n\ne m$ the
inner product of $g\circ F^n$ and $g\circ F^m$ contains the integral
of $\exp(2\pi i(n-m)x)$ over the circle. Its orbit is not precompact.

Form the mapping torus
\[
 N=(\T^2\times\R)/\bigl((z,t)\sim(Fz,t-1)\bigr).
\]
The deck action is free and properly discontinuous, and
$\T^2\times[0,1)$ is a fundamental domain. Hence $N$ is a compact
connected smooth three-manifold. The invariant closed forms
$\dd x\wedge\dd y$ and $\dd t$ descend to forms $\beta$ and $\tau$.
Their product $\nu=\beta\wedge\tau$ is a volume form of total mass one.
The volume-preserving suspension flow is
\[
 \psi_s[z,t]=[z,t+s],\qquad \psi_n[z,t]=[F^n z,t]\quad(n\in\Z).
\]

The flow $\psi$ is ergodic. Indeed, an invariant $v\in L^2(N,\nu)$
has zero distributional derivative in $t$ in each product chart and so is
$v(z,t)=w(z)$ almost everywhere. The seam gives $w\circ F=w$;
the preceding Fourier argument makes $w$ constant. Strong continuity of
the Koopman group follows from uniform continuity for continuous functions
on compact $N$, density of those functions in $L^2$, and unitarity.
In particular invariance at all rational times is enough for this argument.

The fundamental-domain observable
\begin{equation}\label{eq:G}
 G([x,y,t])=\exp(2\pi i y),\qquad 0\le t<1,
\end{equation}
is a valid measurable $L^2$ function; values on the measure-zero seam
do not matter. Its integer translates under $\psi$ are orthonormal.
Thus the continuous-time suspension, although ergodic, is not almost periodic.
Its time-one map need not be ergodic, and we do not use such an assertion.

\section{Hamiltonian construction and derivative bounds}

Let $S^1_E=\R/\Z$ and set
\[
 M=N\times S^1_E,\qquad \omega=\beta+\tau\wedge\dd E.
\]
Here $\dd E$ denotes the global circle one-form, not the differential of
a globally real circle coordinate. Both summands are closed, and
\begin{equation}\label{eq:volume}
 \frac{\omega^2}{2}=\beta\wedge\tau\wedge\dd E
\end{equation}
is nowhere zero and has total integral one. Thus $\omega$ is symplectic.
It is nonexact: its integral over a torus fiber is one.

To define a genuine real-valued Hamiltonian, set
\[
 \eta(r)=\begin{cases}0&r\le0,\\e^{-1/r}&r>0,\end{cases}
 \qquad q(r)=\frac{\eta(r)}{\eta(r)+\eta(1-r)}.
\]
The denominator is everywhere positive; flatness of $\eta$ at zero makes
$q$ smooth, with $q=0$ on $(-\infty,0]$ and $q=1$ on $[1,\infty)$.
On the representative interval $0\le E<1$, define
\[
 \chi(E)=q\bigl(16(E-1/16)\bigr)q\bigl(16(7/16-E)\bigr),
 \qquad h(E)=E\chi(E).
\]
The function $h$ is zero near both endpoints, so its periodic extension is
a globally smooth real function on $S^1_E$. On
$I=(1/8,3/8)$ it equals $E$, with $h'=1$ and $h''=0$.
Define $H(q,E)=h(E)$. Contraction gives
\[
 \iota_{\partial_t}\omega=\dd E,\qquad X_H=h'(E)\partial_t.
\]
The global flow is complete by compactness. On $B=N\times I$ it is
\begin{equation}\label{eq:bandflow}
 \Phi_s(q,E)=(\psi_s(q),E),\qquad \mu(B)=|I|=1/4.
\end{equation}
Since $H=E$ in the band, each $N\times\{E\}$ there is a compact
connected energy component, regular along that component. Its normalized
conditional measure is $\nu$. Other components at the same energy value
outside the band do not affect this claim.

On the cover $\T^2\times\R\times S^1_E$, consider the frame
\[
 V_1=\partial_x-t\partial_y,\qquad V_2=\partial_y,
 \qquad V_3=\partial_t,\qquad V_4=\partial_E.
\]
The derivative of $(z,t,E)\mapsto(Fz,t-1,E)$ sends $V_1(t)$ to
$\partial_x+(1-t)\partial_y=V_1(t-1)$ and fixes the other frame vectors.
The frame is therefore global and smooth. Give $M$ the smooth metric in
which it is orthonormal. On $B$, at arbitrary real time $s$,
\[
 D\Phi_s(V_1)=V_1+sV_2,\qquad D\Phi_s(V_j)=V_j\quad(j=2,3,4),
\]
with the right-hand frame evaluated at the image point. This formula includes
seam crossings and the transverse energy direction. The shear and its inverse
have norm at most $1+|s|$, so for every nonzero tangent vector $v$,
\[
 \frac{\|v\|}{1+|s|}\le\|D\Phi_s v\|\le(1+|s|)\|v\|.
\]
Taking logarithms and dividing by $|s|$ proves that all exponents are zero,
at every band point. Uniform equivalence of smooth metrics on compact $M$
makes the conclusion metric independent.

\section{Every positive-volume invariant region fails almost periodicity}

Let $C\subset B$ be measurable and invariant modulo null sets, with
$\mu(C)>0$. Invariance of its indicator for each rational time, followed by
Fubini and a countable intersection, shows that for almost every $E$ the
section $1_C(\cdot,E)$ is invariant under every rational-time $\psi_s$.
Strong continuity extends this to every real time. Ergodicity of $\psi$
then makes the section indicator either zero or one almost everywhere.
The section mass is measurable in $E$, hence for a measurable $J\subset I$,
\[
 C=N\times J\quad\hbox{modulo null sets},\qquad \mu(C)=|J|>0.
\]
Restrict $G$ from \eqref{eq:G} to $C$. Formula \eqref{eq:iterate} gives
\[
 \langle G\circ\Phi_n,G\circ\Phi_m\rangle_{L^2(C,\mu|_C)}
 =|J|\,\delta_{n,m}.
\]
After multiplication by $|J|^{-1/2}$ these form unit vectors at pairwise
distance $\sqrt2$. Their orbit is not relatively compact. Thus $C$ is
not almost periodic. Taking a single energy fiber instead proves the same
statement for its conditional measure, completing the energy claim.

For clarity, the maximal good region in our convention exists modulo null
sets. Almost periodicity passes to invariant subsets by restriction and to
countable disjoint unions: truncate an $L^2$ function to finitely many
summands, use precompactness on each, and control the invariant tail norm.
General countable unions reduce to disjoint ones by successive differences.
Choose almost-periodic regions whose measures tend to their supremum;
their union is almost periodic and attains that supremum. Adding any other
almost-periodic region changes it only by a null set. This gives a maximal
good region. If it met $B$ in positive measure, restriction to that
intersection would contradict the preceding paragraph. The bad region
has no point in $B$ by the exact derivative bounds. The ugly complement
therefore has measure at least $1/4$, proving Theorem~\ref{thm:main}.

\section{Scope, provenance, and verification}

This is an explicit answer to the printed existential question under the
stated spectral convention. The skew shift, its Fourier mechanism, and the
symplectic mapping-torus construction are classical ingredients. We do not
assert historical first priority for these ingredients or for this consequence.
The original source's subsequent speculation about weak mixing is a distinct
claim: the suspension here has clock eigenfunctions $\exp(2\pi i t)$,
so it is not weakly mixing. Those eigenfunctions do not make its entire
Koopman representation almost periodic. The per-energy result ensures
that the example is not merely an artifact of averaging varying integrable
frequencies over energies.

The supporting package preserves the original candidate proof and exact
verification scripts, the prior independent audit, and new adversarial review
records. The original replay gives 3,248 exact author checks and 2,460 checks
from the prior independent implementation, including negative controls.
These finite checks verify algebra and implementation consistency; they do
not prove smooth flatness, Fourier completeness, irrationality, disintegration,
strong continuity, or infinite-time limits. Those steps are supplied in the
written proof above. Python 3 and its standard library suffice for the original
replay; any additional dependency of a supplemental checker is documented.

\paragraph{AI use and review status.}
AI tools were used extensively in developing and verifying the construction,
checking sources, preparing this note, and conducting adversarial reviews.
AI reviews are fallible and are not human peer review or formal proof
certification. This is an unrefereed preprint; no human peer review is claimed.
The author is responsible for the claims. No journal submission is included.

\begin{thebibliography}{9}
\bibitem{knill2000} O. Knill,
\emph{Some intriguing open problems in Hamiltonian dynamics},
talk, Montana State University, October 19, 2000, item 10.
\url{https://people.math.harvard.edu/~knill/seminars/intr/}.
Primary PDF: \url{https://people.math.harvard.edu/~knill/seminars/bozeman/intr-l.pdf}.
\bibitem{knill1998} O. Knill,
Singular continuous spectrum and quantitative rates of weak mixing,
\emph{Discrete and Continuous Dynamical Systems} \textbf{4} (1998), 33--42.
\href{https://doi.org/10.3934/dcds.1998.4.33}{doi:10.3934/dcds.1998.4.33}.
\bibitem{knill2016} O. Knill, \emph{Differential equations on graphs},
Summer 2016, Section 4.
\url{https://people.math.harvard.edu/~knill/pde/pde.pdf}.
\end{thebibliography}
\end{document}
